% =============================================================================
%  最优传输的定性理论 — Villani《Optimal Transport: Old and New》第 I 部分 (第 4–8 章)
%  模式：被动自学（导师讲义 passive-study）  ·  主题：teal  ·  语言：中文
%  来源：C. Villani, Optimal Transport: Old and New, Springer 2009.
%        Part I "Qualitative description of optimal transport", Chapters 4–8.
%        (PDF pp. 55–210 / 印刷页 pp. 39–198)
%
%  所有示意图用 TikZ（teal 配色）重绘，标注 "据 Villani 图 x.y 重绘"。
%  不使用 beamer overlay 命令 (\pause/\onslide/\only/\uncover)。
%
%  CJK 字体：font-config 自动检测 .fonts/SourceHanSerifSC-Medium.otf（思源宋体）
%  数学字体：KpMath（模板默认）
%
%  本文件为「主控文件」：前置内容在此，各节 \input 自 slides_assets/ot/secN-*.tex
%  组装顺序由 lead-codex 指挥官统一维护；各节文件遵循 STYLE-CONTRACT.md。
% =============================================================================
\documentclass[aspectratio=169,10pt]{beamer}
\usepackage[teal]{template-lib/template-lib}
\uselayout{text}
\uselayout{eq}
\uselayout{twocol}
\uselayout{1img}
\uselayout{2img}
\uselayout{table}

\usetikzlibrary{arrows.meta,calc,positioning,patterns,decorations.pathmorphing,%
                fit,backgrounds,shapes.geometric,shapes.misc,intersections}

\graphicspath{{slides_assets/ot/}{slides_assets/ot/source_figures/}}

% ── 记号缩写（全讲义统一，secN 文件不得重复定义）──────────────────────────────
\newcommand{\R}{\mathbb{R}}
\newcommand{\N}{\mathbb{N}}
\newcommand{\E}{\mathbb{E}}
\newcommand{\Prob}{\mathbb{P}}
\newcommand{\X}{\mathcal{X}}
\newcommand{\Y}{\mathcal{Y}}
\newcommand{\Z}{\mathcal{Z}}
\newcommand{\Pp}{\mathcal{P}}                 % 概率测度集
\newcommand{\Ppp}[1]{\mathcal{P}_{#1}}        % P_p（p 阶矩有限）
\newcommand{\Pac}{\mathcal{P}^{\mathrm{ac}}}  % 绝对连续概率测度
\newcommand{\coup}[2]{\Pi(#1,#2)}             % 耦合（转移计划）集合
\newcommand{\law}{\operatorname{law}}
\newcommand{\Id}{\mathrm{Id}}
\newcommand{\supp}{\operatorname{supp}}
\newcommand{\dom}{\operatorname{dom}}
\newcommand{\pf}[2]{{#1}_{\#}#2}              % 推前 T_# mu
\newcommand{\Wp}[1]{W_{#1}}                   % Wasserstein 距离 W_p
\newcommand{\Wtwo}{W_{2}}
\newcommand{\ctr}[1]{{#1}^{c}}               % c-变换 psi^c
\newcommand{\cctr}[1]{{#1}^{cc}}            % 二次 c-变换
\newcommand{\inner}[2]{\langle #1,\,#2\rangle}
\newcommand{\abs}[1]{\lvert #1\rvert}
\newcommand{\norm}[1]{\lVert #1\rVert}
\newcommand{\dd}{\,\mathrm{d}}
\newcommand{\cost}{c}
\newcommand{\half}{\tfrac{1}{2}}
\newcommand{\eps}{\varepsilon}
\newcommand{\To}{\longrightarrow}
\newcommand{\wkto}{\rightharpoonup}           % 弱收敛
\newcommand{\subjto}{\;\text{s.t.}\;}

% ── 被动自学讲义的结构标记 ────────────────────────────────────────────────────
% 关键术语首次出现
% \TLterm{} 已由模板提供
% 难度星标（amssymb \bigstar），accent 色
\newcommand{\diffI}{\textcolor{TLaccent}{\ensuremath{\bigstar}}}
\newcommand{\diffII}{\textcolor{TLaccent}{\ensuremath{\bigstar\bigstar}}}
\newcommand{\diffIII}{\textcolor{TLaccent}{\ensuremath{\bigstar\bigstar\bigstar}}}
% 习题提示行
\newcommand{\hint}[1]{\emph{提示.}~#1}
% “定义 / 定理 / 命题 / 引理” 行内标签（plain-text-first：默认不进彩框）
\newcommand{\deflab}[1]{\textbf{\textcolor{TLprimaryDark}{#1}}}
% 文献注记标签（用于 References / Bibliographical notes 帧，含英文锚点供校验器识别）
\newcommand{\bibnote}[1]{\emph{文献注记 (bibliographical notes).}~#1}

% ── 本地化“要点”标签（要点。Key Takeaway）─────────────────────────────────────
\renewcommand{\TLtakeaway}[1]{%
  \vspace{0.4em}%
  \begin{tikzpicture}
    \node[fill=TLaccentLight, rounded corners=3pt, inner sep=8pt,
          text width=\linewidth-16pt, align=justify] {%
      \small\textcolor{TLaccent}{\textbf{要点。}} #1%
    };
  \end{tikzpicture}%
  \vspace{0.2em}%
}

\TLsetfoot{最优传输的定性理论 \textperiodcentered{} Villani 第 I 部分（第 4--8 章）}

% =============================================================================
